% Owner: Bach. EDA uses Tuong's reference-aligned processed dataset.
\newcommand{\edageometryfigure}[3]{%
  \begin{figure}[htbp]
    \centering
    \includegraphics[width=0.94\linewidth]{figures/eda/#1}
    \caption{#2}
    \label{#3}
  \end{figure}%
}
\section{Data and Scope}

The descriptive analysis uses the 1,050 observations and five variables
produced by the preprocessing procedure in Chapter~3. The input file is
\texttt{internet\_ads\_processed.csv}. The retained variables are \texttt{Height},
\texttt{Width}, \texttt{Aspect\_Ratio}, \texttt{Local}, and \texttt{IsAds}.
This scope follows the five-variable, mean-imputed branch of the reference
report's preprocessing and descriptive analysis. The original 3,279-row
sample provides the starting population of recorded images, while the
results below describe the reduced processed sample.

Missing height and width values were replaced by their original observed
means, all aspect ratios were recomputed as Width/Height, 15 rows missing
Local were removed, and 2,214 repeated rows on the five selected variables
were excluded. The retained records therefore have no missing values.
Extreme geometric values remain in the sample.

\subsection*{R Implementation and Input}
\begin{lstlisting}[style=edaR]
source("scripts/eda_helpers.R")
data <- read.csv(
  "data/processed/internet_ads_processed.csv",
  check.names = FALSE, na.strings = c("", "NA")
)
geometry <- c("Height", "Width", "Aspect_Ratio")
labels <- c("ad.", "nonad.")
stopifnot(nrow(data) == 1050L, ncol(data) == 5L)
stopifnot(!anyNA(data), all(data$IsAds %in% labels))
\end{lstlisting}

The excerpts use the delivered column names and preserve the trailing
periods in the labels. They share the helper functions used by the R EDA
entry point. The reported calculations and figures were reproduced using
the Python implementation and checked against the delivered processed
CSV. The preprocessing execution log records an R run; the integrated
R EDA implementation has not undergone local runtime validation.

\section{Class Distribution and Missing Data}
\begin{lstlisting}[style=edaR]
class_summary <- summarize_labels(data, "IsAds", labels)
missing_after <- colSums(is.na(data))
print(class_summary)
print(missing_after)
\end{lstlisting}

\texttt{summarize\_labels} counts each label and divides by 1,050.
The missing-value check reports the post-processing missing counts.

\begin{table}[htbp]
\centering
\caption{Class distribution before and after reference-aligned preprocessing.}
\label{tab:eda-classes}
\begin{tabular}{lrrrr}
\toprule
Class & Original count & Original share & Processed count & Processed share\\
\midrule
\texttt{ad.} & 459 & 14.00\% & 93 & 8.86\%\\
\texttt{nonad.} & 2,820 & 86.00\% & 957 & 91.14\%\\
Total & 3,279 & 100.00\% & 1,050 & 100.00\%\\
\bottomrule
\end{tabular}
\end{table}
\edageometryfigure{class_distribution.png}{Class distribution in the reduced processed sample (n=1,050).}{fig:eda-classes-plot}

The processed sample contains about 10.29 non-advertisements for each
advertisement. Predicting only \texttt{nonad.} would produce 91.14\%
accuracy within this sample and zero advertisement recall. The reduction
from 14.00\% to 8.86\% advertisements reflects preprocessing decisions,
particularly deduplication on the selected variables. These proportions
characterize the reduced sample rather than all original images.

\section{Geometric Features by Class}

Mean, sample SD, median, and quartiles summarize each feature by label.
Sample SD uses denominator \(n-1\); quartiles use R's type 7 interpolation.
\begin{lstlisting}[style=edaR]
numeric_summary <- summarize_numeric(
  data, geometry, "IsAds", labels
)
print(numeric_summary[numeric_summary$group != "all", ])
dir.create("output/eda-r/figures", recursive = TRUE,
           showWarnings = FALSE)
for (i in seq_along(geometry)) {
  plot_eda_numeric(
    data, geometry[i], "IsAds", labels,
    sprintf("output/eda-r/figures/numeric_%03d.png", i)
  )
}
\end{lstlisting}

The helper calculates statistics overall and within each label. The plot
function produces a pooled histogram and class-specific boxplots for each
feature. The retained extremes remain visible on their original scales.

\begin{table}[htbp]
\centering
\small
\caption{Geometric features by class in the reduced processed sample.}
\label{tab:eda-geometry}
\begin{tabular}{llrrrrr}
\toprule
Feature & Class & n & Mean & SD & Median & [Q1, Q3]\\
\midrule
Height & \texttt{ad.} & 93 & 68.39 & 56.25 & 60.00 & [40.00, 73.00]\\
Width & \texttt{ad.} & 93 & 221.15 & 149.92 & 155.34 & [100.00, 345.00]\\
Aspect\_Ratio & \texttt{ad.} & 93 & 4.29 & 3.30 & 3.00 & [1.58, 6.97]\\
Height & \texttt{nonad.} & 957 & 82.20 & 65.71 & 64.00 & [35.00, 110.00]\\
Width & \texttt{nonad.} & 957 & 138.24 & 90.13 & 119.00 & [78.00, 180.00]\\
Aspect\_Ratio & \texttt{nonad.} & 957 & 2.73 & 3.61 & 1.50 & [1.00, 3.21]\\
\bottomrule
\end{tabular}
\end{table}
\edageometryfigure{numeric_001.png}{Height histogram and class-specific boxplots after preprocessing.}{fig:eda-1}
\edageometryfigure{numeric_002.png}{Width histogram and class-specific boxplots after preprocessing.}{fig:eda-2}
\edageometryfigure{numeric_003.png}{Aspect ratio histogram and class-specific boxplots after preprocessing.}{fig:eda-3}

Advertisements have a higher median width (155.34 px versus 119 px)
and aspect ratio (3.00 versus 1.50), while their median height is slightly
lower (60 px versus 64 px). Both classes exhibit wide ranges and overlap.
A geometric threshold alone therefore does not establish perfect class
separation. The summaries include imputed dimensions and recomputed
ratios, so their distributions incorporate the effects of preprocessing.

\section{Overall Numerical Summary}
\begin{lstlisting}[style=edaR]
overall_summary <- numeric_summary[
  numeric_summary$group == "all",
  c("variable", "n_valid", "mean", "sd", "min",
    "q1", "median", "q3", "max", "iqr")
]
print(overall_summary)
\end{lstlisting}

The overall rows summarize all 1,050 processed records. Each feature has
the same denominator because missing values were resolved upstream.
Rounding is applied only for presentation.

\begin{table}[htbp]
\centering
\small
\caption{Overall geometric statistics after reference-aligned preprocessing.}
\label{tab:eda-overall}
\begin{tabular}{lrrr}
\toprule
Statistic & Height (px) & Width (px) & Aspect ratio\\
\midrule
n & 1050 & 1050 & 1050\\
Mean & 80.97 & 145.58 & 2.8649\\
Sample SD & 65.02 & 99.65 & 3.6082\\
Minimum & 1.00 & 1.00 & 0.0016\\
Q1 & 35.00 & 81.00 & 1.0000\\
Median & 62.00 & 120.00 & 1.5367\\
Q3 & 108.00 & 189.50 & 3.5724\\
Maximum & 640.00 & 640.00 & 60.0000\\
IQR & 73.00 & 108.50 & 2.5724\\
\bottomrule
\end{tabular}
\end{table}

Height, width, and aspect ratio have overall means of 80.97 px,
145.58 px, and 2.86, respectively. Their medians are 62 px, 120 px,
and 1.54. Higher means than medians and the long upper tails in the
histograms indicate right-skewed distributions. A normal Q--Q plot
provides a further visual comparison with a normal distribution.
\begin{lstlisting}[style=edaR]
par(mfrow = c(1, 2))
for (variable in c("Height", "Width")) {
  qqnorm(data[[variable]], main = variable,
         pch = 16, cex = 0.4)
  qqline(data[[variable]], col = "#a76423")
}
\end{lstlisting}

\texttt{qqnorm} compares ordered feature values with theoretical normal
quantiles; \texttt{qqline} adds a reference line through the quartiles.
Departures from a straight line, particularly in the upper tails, indicate
that a normal approximation to the feature distribution is limited.
The full processed sample is retained for this assessment. A visually
straighter subset obtained through selective filtering would describe a
different sample.
\edageometryfigure{geometry_qq.png}{Normal Q--Q plots for processed Height and Width (n=1,050).}{fig:eda-qq}

\section{Relationships between Geometric Features}

All three pairwise comparisons use the same 1,050 processed records.
Pearson's correlation measures linear association through
\[
 r_{XY}=\frac{\sum_i(x_i-\bar{x})(y_i-\bar{y})}
 {\sqrt{\sum_i(x_i-\bar{x})^2\sum_i(y_i-\bar{y})^2}}.
\]
\begin{lstlisting}[style=edaR]
correlation <- cor(data[, geometry], method = "pearson")
print(correlation)
class_colors <- c("ad." = "#28786d", "nonad." = "#a76423")
pairs(data[, geometry], pch = 16, cex = 0.4,
      col = adjustcolor(class_colors[data$IsAds],
                        alpha.f = 0.35))
\end{lstlisting}

\texttt{cor} calculates the correlation matrix, and \texttt{pairs}
plots feature pairs with class-specific colors. Transparent points reduce
overplotting. The Python figure displays each distinct pair once.
\edageometryfigure{geometry_correlations.png}{Pearson correlations in the reduced processed sample (n=1,050).}{fig:eda-correlations}
\edageometryfigure{geometry_scatterplots.png}{Processed geometric feature pairs colored by class (93 ad. and 957 nonad.).}{fig:eda-scatter}

Height and width have a weak positive association (\(r=0.265\)).
Height and aspect ratio have a negative association (\(r=-0.369\)),
while width and aspect ratio have a positive association
(\(r=0.489\)). These values are consistent with the corresponding
rounded relationships reported in the reference analysis.
Aspect ratio is constructed from width and height; these correlations
partly reflect that dependence. Class overlap and extreme values are
visible in the scatterplots. Correlation alone does not establish the
predictive performance of a feature or a causal relationship.

\section{Extreme Values and Interpretation}

Potential outliers are identified using type 7 quartiles and the fences
\[
 \mathrm{IQR}=Q_3-Q_1,\qquad
 L=Q_1-1.5\,\mathrm{IQR},\qquad U=Q_3+1.5\,\mathrm{IQR}.
\]
\begin{lstlisting}[style=edaR]
iqr_flags <- function(x) {
  q <- quantile(x, c(0.25, 0.75), type = 7,
                names = FALSE)
  lower <- q[1] - 1.5 * diff(q)
  upper <- q[2] + 1.5 * diff(q)
  data.frame(lower_fence = lower, upper_fence = upper,
             n_flagged = sum(x < lower | x > upper))
}
lapply(data[, geometry], iqr_flags)
\end{lstlisting}

The calculation flags values strictly outside the fences. These counts
use pooled thresholds; class-specific boxplots use their own distributions.

\begin{table}[htbp]
\centering
\small
\caption{Potential outliers using pooled 1.5-IQR fences after preprocessing.}
\label{tab:eda-outliers}
\begin{tabular}{lrrrr}
\toprule
Feature & Lower fence & Upper fence & Flagged & Rate\\
\midrule
Height & -74.50 & 217.50 & 38 & 3.62\%\\
Width & -81.75 & 352.25 & 60 & 5.71\%\\
Aspect\_Ratio & -2.86 & 7.43 & 84 & 8.00\%\\
\bottomrule
\end{tabular}
\end{table}

The pooled rule flags 38 heights, 60 widths, and 84 aspect ratios.
The counts can overlap because an observation may be flagged on several
features. For advertisement height, the IQR is now 33 px and the
class-specific rule flags 9 of 93 records. No additional observations
are removed during EDA. Extreme sizes can represent valid image formats;
median and IQR provide complementary summaries to mean and SD.

\section{Binary Feature and Limitations}
\begin{lstlisting}[style=edaR]
local_summary <- summarize_binary(
  data, "Local", "IsAds", labels
)
print(local_summary[local_summary$group != "all", ])
\end{lstlisting}

The helper divides the count of ones by the number of non-missing
values in each class. Local equals one for 66 of 93 advertisements
(70.97\%) and 681 of 957 non-advertisements (71.16\%). The observed
difference is approximately 0.19 percentage points.
\edageometryfigure{binary_001.png}{Local prevalence by class in the reduced processed sample.}{fig:eda-local}

The analysis follows the reference report's mean-imputed preprocessing
branch. Mean imputation introduces repeated central values, and removing
records based on five selected variables can merge distinct images with
different omitted features. The resulting summaries apply to the reduced
sample. The reference report's later median-imputed modeling branch is a
separate dataset. Predictive evaluation requires a separately specified
training/test design and imputation parameters estimated on training data.
No inferential tests or model-performance claims are made in this
chapter; those analyses belong to the statistical inference section.
